% !TEX root = ../main.tex
%-----------------------------------------------------------------------
%  Appendix: running the experiment (the operational how-to).
%  \input by ../main.tex -- not compiled on its own.
%-----------------------------------------------------------------------
%-----------------------------------------------------------------------
\section{Running the experiment}
%-----------------------------------------------------------------------

\begin{callout}
\textbf{Snapshot note (added 2026-08-31).} This appendix records the machine
and toolchain state Part~I actually ran under, and is kept as that record. The
adjoint toolchain has since moved to Tapenade-native \code{TAP\_*} hooks:
stock \code{genmake2}, no post-edited generated files (the
\code{forward\_step\_b.f\_modified\_*} staging described below is gone), the
\code{rawTapenade} control build retired --- raw Tapenade output \emph{is}
the working configuration --- \code{ADJ*} dumps free of the tile-seam
artifact and now including \code{ADJetan}, and \code{useGrdchk = .FALSE.}
as the committed DINO default (this plan's per-member instruction, adopted
repo-wide). The commands themselves are unchanged
(\code{build\_tapAdj.sh}, \code{tools/submit.sh} with \code{IMPACTS\_*}
overrides). See \code{notes/references/tapenade\_hooks/} for the full change
note.

\textbf{Update (2026-09-02).} The adjoint that ran this ensemble --- every
call checkpointed, Tapenade's default --- is no longer the default build.
\code{build\_tapAdj.sh} and \code{submit\_tapAdj.sh} are now symlinks to the
profile-guided \code{-nocheckpoint} pair, whose adjoint is bitwise identical in
\code{fc}, \code{adxx\_*} and \code{ADJ*} and runs a five-year window in
9.6\,h instead of 14.1\,h; the build in the table below lives on as
\code{build\_tapAdj\_ckpAll.sh} / \code{build\_tapAdj\_ckpAll/}. Run
directories now carry the build token, so the ensemble adjoints are
\code{DINO\_1deg\_tapAdj\_ckpAll\_5yr\_M1\_run31040} to
\code{\ldots\_M7\_run31046} on scratch, and
\code{build\_tapAdj\_rawTapenade/} has been deleted. The same eight
adjoints were rerun with the new default on 2026-09-02 as jobs 31060--31067
(\code{DINO\_1deg\_tapAdj\_nocheckpoint\_5yr\_\ldots}): bitwise identical
to 31039--31046 in \code{fc}, \code{adxx\_*} and \code{ADJ*}, the four
blow-ups included, in 9.5--9.75\,h each against 14.0--15.7\,h. The
\code{adjViscBoost} build stays checkpoint-everything on purpose: with the
list it is a different regularisation (run 31056 vs 31025).
\end{callout}

\subsection{Environment}

Verified on 2026-08-21 on the machine the experiment will be run from.

\begin{center}
\begin{tabularx}{\textwidth}{@{}Y{0.45}Y{1.55}@{}}
\toprule
Host / cluster & \code{c1-1}, detected as \code{sverdrup} ---
  \code{machine\_env.sh} falls back to sverdrup because \code{\$NERSC\_HOST} is
  unset\\
Repository & \code{/home/tshahriar/Proj\_ImPACTS}\\
Setup & \code{MITgcm\_c69m/mysetups/DINO\_1deg/}\\
\code{SCRATCH\_ROOT} & \code{/scratch2/tshahriar}\\
\code{MPI\_LAUNCHER} & \code{mpirun -n}\\
\code{SBATCH\_EXTRA} & empty, so \code{tools/submit.sh} is exactly \code{sbatch}
  here\\
\bottomrule
\end{tabularx}
\end{center}

Three dependencies live outside the repository, which is why a fresh clone cannot
build on its own:

\begin{center}
\begin{tabularx}{\textwidth}{@{}Y{0.5}Y{1.5}@{}}
\toprule
MPI optfile & \code{\textasciitilde/tools\_and\_software/crios\_computing/
  computing/optfiles/linux\_amd64\_ifort+mpi\_sverdrup}\\
Tapenade & \code{\textasciitilde/tools\_and\_software/tapenade\_3.16-v2-latest/
  bin/tapenade}, on \code{\$PATH}\\
Java, required by Tapenade & OpenJDK 16.0.1\\
\bottomrule
\end{tabularx}
\end{center}

\code{impacts\_check\_env} passes with no warnings. Run it first if anything
looks wrong:

\begin{center}
\footnotesize
\begin{verbatim}
source tools/machine_env.sh && impacts_check_env
\end{verbatim}
\end{center}

Builds present, all dated 2026-08-18:

\begin{center}
\begin{tabularx}{\textwidth}{@{}Y{1.05}Y{0.75}Y{1.2}@{}}
\toprule
\textbf{Directory} & \textbf{Executable} & \textbf{Use}\\
\midrule
\code{build\_frd/} & \code{mitgcmuv} & the forward leg\\
\code{build\_tapAdj/} & \code{mitgcmuv\_tap\_adj}
  & \textbf{the ensemble adjoint}\\
\code{build\_tapAdj\_adjViscBoost/} & \code{mitgcmuv\_tap\_adj}
  & inflated adjoint viscosity; not this experiment\\
\code{build\_tapAdj\_rawTapenade/} & \code{mitgcmuv\_tap\_adj}
  & control build, uncorrected \code{forward\_step\_b.f}; nothing submits it\\
\bottomrule
\end{tabularx}
\end{center}

\begin{callout}
\textbf{\code{code\_tap/} is compiled as-is} --- since 2026-09-02 no build
script copies anything into it: \code{SIZE.h} is the one decomposition, the
plain builds compile the vendored \code{pkg/autodiff} sources, and the
adjViscBoost build takes its four ASTE-derived shadows from
\code{code\_tap/variants/adjViscBoost/} as a second \code{-mods} directory.

Build directories \emph{symlink} their sources back into \code{code\_tap/}, so
an edit there reaches every build directory. For this ensemble: after any
source edit re-run the build script of the build you submit, and never run
bare \code{make} in an older build directory. To check what a build actually
compiled against, diff its generated \code{.f} file, not its symlinked
\code{.h}.
\end{callout}

Inputs are staged: \code{input\_binaries/} (23 files, 179 MB) and
\code{input\_adj\_binaries/ones\_64b.bin}. Both are gitignored, so they exist on
this machine but not in a fresh clone.

Nothing needs moving. The build, the input binaries and both pickups are already
in place, the latter in the 200-year spin-up run directory on scratch.

\subsection{Generating the diffusivity files}

The reference field is spatially uniform, so each member is a constant. Verify
that before generating anything:

\begin{footnotesize}
\begin{verbatim}
import numpy as np
a = np.fromfile('dino_diffKr.bin', dtype='>f8').reshape(36, 198, 51)
assert np.unique(a).size == 1
print(a.flat[0])          # expect 1.2e-05
\end{verbatim}
\end{footnotesize}

Then generate the seven perturbed fields:

\begin{footnotesize}
\begin{verbatim}
import numpy as np

NZ, NY, NX = 36, 198, 51
REF = 1.2e-5

members = {
    'M1': 0.25,   'M2': 0.5,   'M3': 2,   'M4': 4,
    'M5': 8,      'M6': 16,    'M7': 32,
}

for tag, factor in members.items():
    kappa = REF * factor
    np.full((NZ, NY, NX), kappa, dtype='>f8').tofile(f'dino_diffKr_{tag}.bin')
    print(f'{tag}: kappa = {kappa:.2e} m2/s')
\end{verbatim}
\end{footnotesize}

Every file must be $2{,}908{,}224$ bytes. A wrong size means a wrong dtype or
shape; \code{>f8}, big-endian float64, is required because
\code{readBinaryPrec=64}. No control binary is needed, since the completed
reference run is the control.

\subsection{Namelists}
\label{app:namelists}

Namelist variants live in \code{input*/variants/}, \textbf{not} directly in
\code{input*/} --- anything placed directly there is staged into every run.
This ensemble's members are grouped in their own directory, so the submit
script resolves \code{test\_cases="kappa\_v\_ensemble/M<n>"} to
\code{variants/kappa\_v\_ensemble/data\_M<n>} and stages only that one. A tag
containing \code{/} names one member of an experiment directory; a bare tag
still names a flat \code{variants/data\_<tag>}.

Forward namelist, \code{input/variants/kappa\_v\_ensemble/data\_M<n>}, for the
leg from 2170 to 2180:

\begin{footnotesize}
\begin{verbatim}
nIter0      = 2986560,       <- nearest checkpoint below 2170
nTimeSteps  = 175680,        <- = 3162240 - nIter0, lands exactly on 2180
diffKrFile  = 'dino_diffKr_M<n>.bin',
\end{verbatim}
\end{footnotesize}

Adjoint namelist, \code{input\_tap/variants/kappa\_v\_ensemble/data\_M<n>},
identical to the reference adjoint apart from the diffusivity file:

\begin{footnotesize}
\begin{verbatim}
nIter0      = 3162240,       <- the pickup the forward leg writes; 2180
nTimeSteps  = 87840,         <- 5 years, 2180 -> 2185
diffKrFile  = 'dino_diffKr_M<n>.bin',
\end{verbatim}
\end{footnotesize}

Everything else stays identical to the reference, since the entire design rests
on $\kv$ being the only difference. The timing parameters in particular are the
same numbers.

\paragraph{Three settings to fix before submitting}, all cheap now and expensive
later:

\begin{enumerate}
  \item \textbf{\code{useGrdchk = .FALSE.}} in \code{data.pkg} for every member.
        This is the largest single saving available
        (Section~\ref{sec:grdchk-cost}), and nothing is lost because the check
        cannot pass at its current point and Section~\ref{sec:fd} is the real
        validation.
  \item \textbf{\code{pChkptFreq}.} The current value writes a permanent restart
        every 30 model days: 120 files and about 2.8 GB per member, of which only
        the last is needed. Raising it saves roughly 20 GB across the ensemble.
        It must \textbf{not} be raised for the spin-up, whose $2{,}400$ retained
        restarts are the Axis~2 dataset.
  \item \textbf{Forward diagnostic cadence.} If the surrogate proves to be a
        local operator (open question 5) it needs the forward state at the same
        cadence as the sensitivities. Forward diagnostics are currently monthly
        while adjoint output is 5-daily. Matching them means editing
        \code{data.diagnostics} before the runs; discovering the need afterwards
        means re-running the ensemble.
\end{enumerate}

\subsection{Submission}

No script changes are required. The submit scripts select a namelist via
\code{test\_cases}, and everything machine-specific comes from
\code{tools/machine\_env.sh}. Per member, set
\code{IMPACTS\_TEST\_CASE=kappa\_v\_ensemble/M<n>} in the environment, which
leaves the working tree clean --- and edit the pickup \code{ln -s} further
down, which is not auto-patched --- then:

\begin{footnotesize}
\begin{verbatim}
cd MITgcm_c69m/mysetups/DINO_1deg
../../../tools/submit.sh submit_frd.sh      # the 2170 -> 2180 leg
../../../tools/submit.sh submit_tapAdj.sh   # the 2180 -> 2185 adjoint
\end{verbatim}
\end{footnotesize}

Use \code{tools/submit.sh} rather than bare \code{sbatch}: on this machine they
are identical, but on Perlmutter bare \code{sbatch} silently omits the account
and QOS.

\begin{callout}
\textbf{\code{nIter0} is not one of the parameters the submit script patches.}
It is baked into whichever \code{data\_<tag>} \code{test\_cases} selects, while
the pickup is a hardcoded \code{ln -s} further down the same script. Changing the
duration is safe; changing the starting point means editing both, and a member
whose \code{nIter0} and pickup symlink disagree will simply fail to find its
pickup. Every member in this ensemble changes the starting point.
\end{callout}

\begin{callout}
\textbf{Submitting edits the tracked namelist.} The submit scripts apply
\code{sed -i} to the duration parameters directly in
\code{\$SLURM\_SUBMIT\_DIR/input\_tap/}, so running a member dirties the working
tree. Expect the diff and inspect it before committing; \code{git status} alone
will not distinguish a change the user authored from one a submit script made.
\code{./tools/pre\_push\_check.sh} flags this specifically.
\end{callout}

\subsection{Running the pilot}

Run the pilot on this machine, since nothing needs porting or transferring. Its
output then becomes the verification target for any later port.

\begin{enumerate}
  \item \textbf{Forward leg}, 2170 $\rightarrow$ 2180, about 1.6 hours at the
        measured rate.
  \item \textbf{Check re-equilibration.} Plot overturning strength over the ten
        years. If it is still drifting at year 10, the sensitivities reflect
        drift rather than a new equilibrium; this is open question 4.
  \item \textbf{Check convection.} \code{ivdc\_kappa = 100} applies a very large
        implicit diffusivity wherever the column is unstable. At 8 $\times$
        $\kv$, explicit mixing may suppress convection, in which case part of the
        sensitivity change is convection triggering less often rather than a
        smooth response to $\kv$. Diagnose convective activity so the two can be
        separated.
  \item \textbf{Adjoint, 5 years}, about 9 hours with the gradient check
        disabled, 23 hours with it enabled.
  \item \textbf{Compare against the reference run} using the same fields, level
        and colour scale as the existing animations. Both cover
        2180 $\rightarrow$ 2185, so any difference is attributable to $\kv$.
\end{enumerate}

Then review before spending the remainder.

\subsection{Remaining members}

The six remaining members are independent and can be submitted together, four per
node if the benchmark shows packing is efficient. Each follows the same sequence:
forward leg, new pickup, adjoint. The total is roughly 64 node-hours with the
gradient check disabled.

The strongly perturbed members, M6 and M7, may alter stratification enough to
require a shorter timestep or more solver iterations. If a member goes unstable,
reduce \code{deltaT} for that member alone and record it: doing so breaks the
``only $\kv$ differs'' property and must be noted in the analysis.

\subsection{Porting to another machine}

Should the ensemble move to Perlmutter, most of the work is already done. Machine
differences were factored out of the scripts, so porting is configuration rather
than editing; \code{PORTING.md} in the repository root is the walkthrough.

\begin{center}
\begin{tabularx}{\textwidth}{@{}Y{0.6}Y{1.4}@{}}
\toprule
\textbf{Piece} & \textbf{Where}\\
\midrule
Per-machine settings & \code{tools/machine\_env.sh} --- detects Perlmutter from
  \code{\$NERSC\_HOST} and sets \code{SCRATCH\_ROOT}, \code{MPI\_LAUNCHER}
  (\code{srun -n}), both optfiles and \code{SBATCH\_EXTRA}\\
Perlmutter optfile & \code{tools/optfile\_templates/linux\_amd64\_gnu+mpi\_%
  perlmutter} --- written but \textbf{not yet validated on Perlmutter}\\
Account, QOS, constraint, walltime & passed by \code{tools/submit.sh} on the
  sbatch command line\\
Environment check & \code{impacts\_check\_env} warns about a missing Tapenade
  binary, missing \code{input\_binaries/}, and an unset \code{NERSC\_ACCOUNT}\\
\bottomrule
\end{tabularx}
\end{center}

What must be transferred, since it cannot be regenerated:

\begin{center}
\begin{tabularx}{\textwidth}{@{}Y{1.15}Y{1.55}Y{0.3}@{}}
\toprule
\textbf{What} & \textbf{Source} & \textbf{Size}\\
\midrule
\code{input\_binaries/} & outside the repo; untracked & 179 MB\\
\code{input\_adj\_binaries/ones\_64b.bin} & outside the repo; small but
  \textbf{not optional} & small\\
\code{pickup.0003162240.\{data,meta\}} (year 2180) & this machine's scratch, run
  28463 & 20 MB\\
\code{pickup.0002986560.\{data,meta\}} ($\approx$2170) & same directory & 20 MB\\
\bottomrule
\end{tabularx}
\end{center}

\textbf{The two pickups are the critical items.} Without them every member needs
its own multi-century spin-up, which is two orders of magnitude more compute.
Move them first, keep second copies, and use Globus.

What then remains, in order of risk:

\begin{enumerate}
  \item \textbf{Tapenade requires a Java runtime at build time}, and Tapenade is
        not in the repository --- \code{genmake2 -tap} calls a \code{tapenade}
        binary on \code{\$PATH}. Confirm this first; it is the most likely thing
        to block the build outright.
  \item \textbf{Validate the optfile.} It has never been compiled on the target.
        Compiler wrappers are \code{ftn} and \code{cc}, not \code{mpif77} and
        \code{mpicc}.
  \item \textbf{Set \code{NERSC\_ACCOUNT}} to the project allocation.
  \item \textbf{Walltime.} The submit scripts carry \code{-t 240:00:00}, which
        exceeds every Perlmutter QOS. Set it from the benchmark: a member needs
        about 11 hours.
  \item \textbf{Scratch purging.} NERSC purges unaccessed scratch files, so
        anything worth keeping must move to the community filesystem.
\end{enumerate}

\code{run\_timing.txt} is written into every run directory by the submit scripts.
It should be kept: it is what turned this project's cost from a guess into the
measurements in Section~\ref{sec:compute}.

